% =========================================================
% Non-Euclidean Geometry — Notes Index
% =========================================================

% ---------------------------------------------------------
% Breadcrumb
% ---------------------------------------------------------

% ---------------------------------------------------------
% Structural Role
% ---------------------------------------------------------
\begin{tcolorbox}[
  colback=gray!6,
  colframe=gray!40,
  arc=2pt,
  left=8pt, right=8pt, top=6pt, bottom=6pt,
  title={\small\textbf{Structural Role: Non-Euclidean Geometry — Geometry is a Choice}},
  fonttitle=\small\bfseries
]
Non-Euclidean geometry is the realisation that \emph{geometry
is a choice of axioms}, not an inevitable description of space.

The Poincar\'e disk is a complete, consistent model of hyperbolic
geometry living entirely within Euclidean geometry.
Its existence proves that the parallel postulate cannot be
derived from the other four --- if it could, hyperbolic geometry
would be inconsistent.
This is model-theoretic reasoning in action: independence is
established by exhibiting a model.

The curvature interpretation comes later, in Riemannian geometry,
where each non-Euclidean geometry is realised as a space of
constant sectional curvature.
The present chapter builds the geometric intuition; Riemannian
geometry provides the analytic framework.
\end{tcolorbox}
\vspace{1em}

% ---------------------------------------------------------
% Structural Roadmap
% ---------------------------------------------------------
\subsection*{Structural Roadmap}

\begin{center}
\textbf{Definitions $\longrightarrow$ Main Theorems
$\longrightarrow$ Consequences and Structural Insight}
\end{center}

The global progression is:

\begin{enumerate}
  \item The parallel postulate and its historical status
  \item Spherical geometry: great circles, angle sums, area
  \item The Beltrami--Klein model of hyperbolic geometry
  \item The Poincar\'e disk and upper half-plane models
  \item Hyperbolic isometries and M\"obius transformations
  \item Angle sums and the defect formula in hyperbolic geometry
  \item The Gauss--Bonnet theorem in the non-Euclidean setting
  \item Elliptic geometry and projective models
  \item Constant-curvature spaces as Riemannian manifolds
\end{enumerate}

\begin{remark*}[Primary sources]
Greenberg, \textit{Euclidean and Non-Euclidean Geometries};
Stillwell, \textit{Sources of Hyperbolic Geometry};
Ratcliffe, \textit{Foundations of Hyperbolic Manifolds}.
\end{remark*}

% ---------------------------------------------------------
% Status
% ---------------------------------------------------------
\vspace{1em}
\begin{tcolorbox}[
  colback=gray!6, colframe=gray!40, arc=2pt,
  left=8pt, right=8pt, top=6pt, bottom=6pt,
  title={\small\textbf{Status: Planned}},
  fonttitle=\small\bfseries
]
\begin{center}\Large\bfseries Coming Soon\end{center}
\vspace{6pt}
\noindent Notes, proofs, and exercises will appear here in a future revision.
\end{tcolorbox}

% Future sections (uncomment as content is written):
% \input{volume-v/geometry/non-euclidean-geometry/notes/notes-parallel-postulate}
% \input{volume-v/geometry/non-euclidean-geometry/notes/notes-spherical}
% \input{volume-v/geometry/non-euclidean-geometry/notes/notes-klein-model}
% \input{volume-v/geometry/non-euclidean-geometry/notes/notes-poincare-disk}
% \input{volume-v/geometry/non-euclidean-geometry/notes/notes-isometries}
% \input{volume-v/geometry/non-euclidean-geometry/notes/notes-gauss-bonnet}
% \input{volume-v/geometry/non-euclidean-geometry/notes/notes-constant-curvature}
